\section*{Bab \ref{ch:b1:fractions} --- Pecahan Rasional}

\begin{solution}{exo:b1:fractions:1}
$\dfrac{1}{X^2 - 1}$: dengan \omterm{def:b1:fractions:field}{kutub} sederhana $\pm 1$; lalu penutupan memberikan
$\dfrac{1/2}{X-1} - \dfrac{1/2}{X+1}$.

$\dfrac{X}{X^2 - 3X + 2} = \dfrac{X}{(X-1)(X-2)}$: penutupan memberikan
$\frac{1}{1-2} = -1$ di $1$ dan $\frac{2}{2-1} = 2$ di $2$, sehingga
$\dfrac{-1}{X-1} + \dfrac{2}{X-2}$.

$\dfrac{X^2+1}{X(X-1)}$: derajatnya $0$, jadi ada bagian
bulatnya: dengan membaginya, $X^2 + 1 = (X^2 - X) + (X + 1)$, sehingga $F = 1 +
\frac{X+1}{X(X-1)}$. Penutupan pada sisanya: $\frac{1}{-1} = -1$
di $0$, dan $\frac{2}{1} = 2$ di $1$:
\[
F = 1 - \frac{1}{X} + \frac{2}{X-1} .
\]
\end{solution}

\begin{solution}{exo:b1:fractions:2}
$\dfrac{1}{X(X^2+1)}$: dengan bentuk $\frac aX + \frac{bX + c}{X^2 + 1}$.
Penutupan di $0$: $a = 1$. Limit $xF$: $0 = a + b$, jadi $b = -1$.
Perhitungan nilai di $X = 1$: $\frac12 = 1 + \frac{c - 1}{2}$, jadi $c = 0$:
\[
\frac{1}{X(X^2+1)} = \frac 1X - \frac{X}{X^2+1} .
\]

$\dfrac{X^3}{X^2+X+1}$: dengan pembagian: $X^3 = (X^2+X+1)(X - 1) + 1$, jadi
\[
\frac{X^3}{X^2+X+1} = X - 1 + \frac{1}{X^2 + X + 1} ,
\]
yang sudah berbentuk terurai secara real (karena kuadratnya berdiskriminan
negatif).
\end{solution}

\begin{solution}{exo:b1:fractions:3}
$\dfrac{1}{X^2(X-1)}$: dengan bentuk $\frac{a}{X^2} + \frac bX +
\frac{c}{X-1}$. Penutupan pada \omterm{def:b1:fractions:field}{kutub} ganda $0$:
$a = \bigl[\frac{1}{X-1}\bigr]_{0} = -1$. Penutupan di $1$: $c = 1$.
Limit $xF$: $0 = b + c$, jadi $b = -1$:
\[
\frac{1}{X^2(X-1)} = -\frac{1}{X^2} - \frac1X + \frac{1}{X-1} .
\]

$\dfrac{X+1}{(X-1)^3}$: dengan $Y = X - 1$, pembilangnya menjadi $Y + 2$:
\[
\frac{Y + 2}{Y^3} = \frac{1}{Y^2} + \frac{2}{Y^3}
= \frac{1}{(X-1)^2} + \frac{2}{(X-1)^3} .
\]
\end{solution}

\begin{solution}{exo:b1:fractions:4}
\omterm{def:b1:fractions:field}{Kutubnya} adalah \omterm{def:b1:complex:unity}{akar satuan} keempat $1, \iu, -1, -\iu$, yang semuanya
sederhana. Rumus \omterm{def:b1:fractions:field}{kutub} sederhananya dengan $B' = 4X^3$: koefisien di $a$ adalah
$\frac{1}{4a^3} = \frac{a}{4a^4} = \frac a4$ (memakai $a^4 = 1$). Jadi,
atas $\C$:
\[
\frac{1}{X^4 - 1}
= \frac{1/4}{X - 1} - \frac{1/4}{X + 1}
+ \frac{\iu/4}{X - \iu} - \frac{\iu/4}{X + \iu} .
\]
Dengan mengelompokkan pasangan \omterm{def:b1:complex:field}{konjugatnya} (yang berpenyebut persekutuan $X^2 + 1$):
$\frac{\iu}{4}\bigl(\frac{1}{X-\iu} - \frac{1}{X+\iu}\bigr) =
\frac{\iu}{4}\cdot\frac{2\iu}{X^2+1} = \frac{-1/2}{X^2+1}$. Jadi atas
$\R$:
\[
\frac{1}{X^4 - 1}
= \frac{1/4}{X-1} - \frac{1/4}{X+1} - \frac{1/2}{X^2 + 1} .
\]
\emph{Periksa di $X = 0$:} $-1 = -\frac14 - \frac14 - \frac12$.
\end{solution}

\begin{solution}{exo:b1:fractions:5}
$\dfrac{X^2}{(X^2+1)^2} = \dfrac{(X^2 + 1) - 1}{(X^2+1)^2} =
\dfrac{1}{X^2+1} - \dfrac{1}{(X^2+1)^2}$.

Jadi sebuah antiturunannya:
\[
\int \frac{x^2\,\dd x}{(x^2+1)^2}
= \arctan x - \frac12\Bigl(\arctan x + \frac{x}{x^2+1}\Bigr) + C
= \frac12\arctan x - \frac{x}{2(x^2+1)} + C .
\]
\end{solution}

\begin{solution}{exo:b1:fractions:6}
\omterm{def:b1:fractions:field}{Kutubnya} $0, -1, \dots, -n$ semuanya sederhana. Penutupan di $-k$:
\[
c_k = \frac{n!}{\prod_{j \neq k} (-k + j)}
= \frac{n!}{\bigl(\prod_{j=0}^{k-1}(j - k)\bigr)
\bigl(\prod_{j=k+1}^{n}(j-k)\bigr)}
= \frac{n!}{(-1)^k k!\,(n-k)!} = (-1)^k \binom nk .
\]
Jadi
\[
\frac{n!}{X(X+1)\cdots(X+n)}
= \sum_{k=0}^{n} \frac{(-1)^k \binom nk}{X + k} .
\]
(Pemeriksaan kewajaran untuk $n = 1$: $\frac{1}{X(X+1)} = \frac1X -
\frac{1}{X+1}$.)
\end{solution}

\begin{solution}{exo:b1:fractions:7}
$P = X^n - 1$ mempunyai $n$ akar sederhana $\omega^k$
(\cref{thm:b1:complex:roots}), jadi kesamaan turunan logaritmik
pada \cref{exo:b1:poly:11} berbunyi
\[
\sum_{k=0}^{n-1} \frac{1}{X - \omega^k}
= \frac{P'(X)}{P(X)} = \frac{n X^{n-1}}{X^n - 1} .
\]
Di $X = 2$, $n = 4$: ruas kanannya $= \frac{4 \times 8}{15} =
\frac{32}{15}$. Ruas kirinya: $\frac{1}{2-1} + \frac{1}{2+1} +
\frac{1}{2 - \iu} + \frac{1}{2 + \iu} = 1 + \frac13 + \frac{4}{5} =
\frac{15 + 5 + 12}{15} = \frac{32}{15}$, dengan memakai $\frac{1}{2-\iu} +
\frac{1}{2+\iu} = \frac{4}{5}$.
\end{solution}

\begin{solution}{exo:b1:fractions:8}
Penutupan: $\dfrac{1}{k(k+1)(k+2)} = \dfrac{1/2}{k} - \dfrac{1}{k+1} +
\dfrac{1/2}{k+2}$. Tulis ulang sebagai selisih yang teleskopik:
\[
\frac{1}{k(k+1)(k+2)}
= \frac12\Bigl(\frac{1}{k(k+1)} - \frac{1}{(k+1)(k+2)}\Bigr),
\]
(jabarkan untuk memeriksanya --- atau kurangkan kedua penguraiannya). Dengan menjumlahkannya:
\[
S_n = \frac12\Bigl(\frac{1}{1 \times 2} - \frac{1}{(n+1)(n+2)}\Bigr)
= \frac14 - \frac{1}{2(n+1)(n+2)}
\xrightarrow[n \to \infty]{} \frac14 .
\]
\end{solution}

\begin{solution}{exo:b1:fractions:9}
Uraikan, untuk $0 \leq m \leq n - 1$, pecahan $\frac{X^m}{P}$
(yang berderajat $m - n \leq -1$, dengan \omterm{def:b1:fractions:field}{kutub} sederhana): maka rumus \omterm{def:b1:fractions:field}{kutub} sederhananya memberikan
\[
\frac{X^m}{P} = \sum_{i=1}^{n} \frac{x_i^m / P'(x_i)}{X - x_i} .
\]
Kalikan dengan $X$ lalu ambil $X \to +\infty$: ruas kirinya menuju
limit $X^{m+1}/P$, yang bernilai $0$ bila $m \leq n - 2$ dan $1$ bila $m =
n - 1$ (karena $P$ \omterm{def:b1:poly:def}{monik} berderajat $n$); sedangkan ruas kanannya menuju
$\sum_i \frac{x_i^m}{P'(x_i)}$. Jadi
\[
\sum_{i=1}^{n} \frac{x_i^{m}}{P'(x_i)} =
\begin{cases}
0 & \text{untuk } 0 \leq m \leq n-2,\\
1 & \text{untuk } m = n - 1,
\end{cases}
\]
yang memuat kedua kesamaan yang dinyatakan tadi (dengan $m = 0$ menuntut $n \geq
2$). (Tafsirannya lewat \cref{thm:b1:poly:lagrange}: jumlah itu adalah
koefisien utama penginterpolasi Lagrange bagi $X^m$, dan
menginterpolasi \omterm{def:b1:poly:def}{polinomial} berderajat $\leq n-1$ pada $n$ titik
menghasilkannya kembali secara persis.)
\end{solution}

\begin{solution}{exo:b1:fractions:10}
Bentuknya $\frac aX + \frac b{X+1} + \frac c{(X+1)^2}$. Penutupan di
$0$: $a = 1$; penutupan pada \omterm{def:b1:fractions:field}{kutub} gandanya: $c = \bigl[\frac1X
\bigr]_{X=-1} = -1$; dan limit $xF(x)$ di tak hingga: $0 = a + b$, jadi
$b = -1$:
\[
\frac{1}{X(X+1)^2} = \frac1X - \frac1{X+1} - \frac1{(X+1)^2} .
\]
Dengan menjumlahkannya untuk $n = 1, \dots, N$: kedua bata pertamanya berteleskop menjadi
$1 - \frac1{N+1}$, dan yang ketiga menyumbang $-\sum_{k=2}^{N+1}
\frac1{k^2}$. Lalu dengan mengambil $N \to \infty$ dan memakai $\sum_{k\geq1}
\frac1{k^2} = \frac{\pi^2}6$:
\[
\sum_{n=1}^{\infty}\frac{1}{n(n+1)^2}
= 1 - \Bigl(\frac{\pi^2}6 - 1\Bigr) = 2 - \frac{\pi^2}6
\approx 0.355 .
\]
\end{solution}

\begin{solution}{exo:b1:fractions:11}
Dalam variabel $Y$: $\frac1{(Y+1)(Y+4)} = \frac{1/3}{Y+1} -
\frac{1/3}{Y+4}$ (dari penutupan di $-1$ dan $-4$), jadi
\[
\frac{1}{(X^2+1)(X^2+4)}
= \frac13\,\frac1{X^2+1} - \frac13\,\frac1{X^2+4} .
\]
Demikian pula $\frac{Y}{(Y+1)(Y+4)} = \frac{-1/3}{Y+1} +
\frac{4/3}{Y+4}$, sehingga
\[
\frac{X^2}{(X^2+1)(X^2+4)}
= -\frac13\,\frac1{X^2+1} + \frac43\,\frac1{X^2+4} .
\]
(Periksa di $X = 0$: $0 = \frac13(-1 + 1)$.) Keduanya sudah merupakan
penguraian realnya: karena pembilang di atas kuadrat yang tak tereduksikan
kebetulan berupa konstanta.
\end{solution}

\begin{solution}{exo:b1:fractions:12}
\begin{enumerate}
  \item Pada setiap selang yang menghindari \omterm{def:b1:fractions:field}{kutubnya}, $F'(x) = -\sum_i
        \frac{c_i}{(x - p_i)^2} < 0$: jadi ia turun tegas. Ketika $x
        \to \pm\infty$, setiap batanya menuju $0$: sehingga $F \to 0$, dari
        atas di $+\infty$ (karena semua batanya positif di sana) dan dari
        bawah di $-\infty$. Ketika $x \to p_i^+$, bata
        $\frac{c_i}{x - p_i}$ meledak ke $+\infty$ sedangkan yang lain
        tetap terbatas: jadi $F \to +\infty$; demikian pula $F \to -\infty$ ketika
        $x \to p_i^-$.
  \item Tetapkan $\lambda > 0$. Pada $\intoo{-\infty}{p_1}$: $F$
        turun dari $0^-$ ke $-\infty$, jadi $F < 0 < \lambda$:
        tak ada penyelesaian. Pada masing-masing $\intoo{p_i}{p_{i+1}}$ ($1 \leq i
        \leq r-1$): $F$ turun dari $+\infty$ ke $-\infty$,
        sehingga ia bernilai $\lambda$ tepat sekali (menurut sifat nilai
        antara beserta kemonotonan tegasnya). Pada
        $\intoo{p_r}{+\infty}$: $F$ turun dari $+\infty$ ke
        $0^+$, sehingga sekali lagi tepat satu penyelesaian. Totalnya: tepat $r$
        penyelesaian, yang berselang-seling dengan \omterm{def:b1:fractions:field}{kutubnya}.
\end{enumerate}
\end{solution}

\begin{solution}{pb:b1:fractions:1}
\textbf{1.} Penutupan: $\frac1{X(X+1)} = \frac1X - \frac1{X+1}$.
Jumlahnya pun berteleskop:
\[
\sum_{n=1}^{N}\Bigl(\frac1n - \frac1{n+1}\Bigr)
= 1 - \frac1{N+1} \longrightarrow 1 .
\]

\textbf{2.} $\frac1{X(X+2)} = \frac{1/2}X - \frac{1/2}{X+2}$.
Dengan menjumlahkannya, suku $\frac1n$ bertahan untuk $n = 1, 2$ dan suku
$-\frac1{n+2}$ bertahan untuk $n = N-1, N$:
\[
\sum_{n=1}^{N}\frac1{n(n+2)}
= \frac12\Bigl(1 + \frac12 - \frac1{N+1} - \frac1{N+2}\Bigr)
\longrightarrow \frac34 .
\]

\textbf{3.} Jika $F(X) = G(X) - G(X+1)$, maka $\sum_{n=1}^N F(n) =
\sum_{n=1}^N\bigl(G(n) - G(n+1)\bigr) = G(1) - G(N+1)$: karena semua
nilai antaranya saling hapus berpasangan. Untuk $F =
\frac1{X(X+1)(X+2)}$, saksinya adalah $G(X) = \frac1{2X(X+1)}$:
\[
G(X) - G(X+1)
= \frac{(X+2) - X}{2X(X+1)(X+2)} = F(X) ,
\]
jadi $\sum_{n=1}^N F(n) = \frac14 - \frac1{2(N+1)(N+2)} \to
\frac14$, yaitu nilai pada \cref{exo:b1:fractions:8}.

\textbf{4.} Letakkan ruas kanannya di atas penyebut persekutuan
$X(X+1)\cdots(X+k)$:
\[
\frac1k\cdot\frac{(X + k) - X}{X(X+1)\cdots(X+k)}
= \frac{1}{X(X+1)\cdots(X+k)} ,
\]
yang tepat merupakan kesamaannya. Jadi $F_k(X) = \frac1k\bigl(G_k(X) -
G_k(X+1)\bigr)$ dengan $G_k(X) = \frac1{X(X+1)\cdots(X+k-1)}$, dan
mekanisme pertanyaan 3 memberikan
\[
\sum_{n=1}^{N}\frac1{n(n+1)\cdots(n+k)}
= \frac1k\Bigl(\frac1{k!} - G_k(N+1)\Bigr)
\longrightarrow \frac1{k\cdot k!} ,
\]
karena $G_k(1) = \frac1{k!}$ dan $G_k(N+1) \to 0$. Untuk $k = 2$:
$\frac1{2\cdot2} = \frac14$, yang cocok dengan pertanyaan 3.

\textbf{5.} \cref{exo:b1:fractions:6} memberikan
$\frac{n!}{X(X+1)\cdots(X+n)} = \sum_{j=0}^n
\frac{(-1)^j\binom nj}{X + j}$. Hitung nilainya di $X = 1$: ruas kirinya
menjadi $\frac{n!}{(n+1)!} = \frac1{n+1}$, sedangkan ruas kanannya menjadi
$\sum_j(-1)^j\binom nj\frac1{1+j}$: itulah kesamaan pertamanya. Di $X = 2$:
ruas kirinya menjadi $\frac{n!}{2\cdot3\cdots(n+2)} =
\frac{n!\cdot1}{(n+2)!} = \frac1{(n+1)(n+2)}$, sedangkan ruas kanannya
$\sum_j(-1)^j\binom nj\frac1{2+j}$: itulah kesamaan keduanya.

\textbf{6.} \omterm{def:b1:fractions:field}{Kutub} $\omega^k$ semuanya sederhana, dan rumus \omterm{def:b1:fractions:field}{kutub}
sederhana pada \cref{met:b1:fractions:compute} memberikan koefisien
\[
\frac{1}{\bigl(nX^{n-1}\bigr)_{X = \omega^k}}
= \frac{1}{n\,\omega^{k(n-1)}}
= \frac{\omega^k}{n\,\omega^{kn}} = \frac{\omega^k}n ,
\]
dengan memakai $\omega^{kn} = 1$. Jadi $\frac1{X^n-1} = \frac1n\sum_k
\frac{\omega^k}{X - \omega^k}$.

\textbf{7.} Untuk $n = 2$ (dengan $\omega = -1$): $\frac12\bigl(
\frac1{X-1} - \frac1{X+1}\bigr) = \frac12\cdot\frac{2}{X^2-1} =
\frac1{X^2-1}$: jadi benar. Koefisiennya berjumlah
$\frac1n\sum_k\omega^k = 0$ untuk $n \geq 2$
(\cref{prop:b1:complex:sumroots}). Memang harus demikian: karena $x\,F(x) \to
\sum_k c_k$ ketika $x \to \infty$ untuk sebarang penguraian berkutub
sederhana, sedangkan di sini $xF(x) = \frac{x}{x^n-1} \to 0$ karena $n \geq
2$.

\textbf{8.} Penutupan untuk $\frac{X^{n-1}}{X^n - 1}$ di $\omega^k$:
$\frac{A(\omega^k)}{B'(\omega^k)} =
\frac{\omega^{k(n-1)}}{n\omega^{k(n-1)}} = \frac1n$, jadi
$\frac{X^{n-1}}{X^n-1} = \frac1n\sum_k\frac1{X - \omega^k}$ ---
yaitu kesamaan \cref{exo:b1:fractions:7} sekali lagi.

\textbf{9.} Dengan $c = \omega^k$, $\conj c = \omega^{n-k}$ dan
$c\conj c = 1$, $c + \conj c = 2\cos\theta_k$:
\[
\frac{c}{X - c} + \frac{\conj c}{X - \conj c}
= \frac{c(X - \conj c) + \conj c(X - c)}
{(X - c)(X - \conj c)}
= \frac{2\cos\theta_k\,X - 2}{X^2 - 2\cos\theta_k\,X + 1} .
\]
Dengan mengelompokkan $k$ bersama $n - k$ pada pertanyaan 6: untuk $n$ ganjil,
\[
\frac1{X^n - 1} = \frac1n\Biggl(\frac1{X-1} +
\sum_{k=1}^{(n-1)/2}
\frac{2\cos\theta_k X - 2}{X^2 - 2\cos\theta_k X + 1}\Biggr) ;
\]
sedangkan untuk $n$ genap, \omterm{def:b1:fractions:field}{kutub} $\omega^{n/2} = -1$ yang berpasangan dengan
dirinya sendiri menyumbang $\frac{-1}{X+1}$ di dalam kurungnya dan jumlah
pasangannya berjalan sampai $\frac n2 - 1$.

\textbf{10.} Untuk $n = 4$: $\theta_1 = \frac\pi2$, $\cos\theta_1 = 0$,
jadi suku pasangannya adalah $\frac{-2}{X^2+1}$ dan
\[
\frac1{X^4-1} = \frac14\Bigl(\frac1{X-1} - \frac1{X+1}
- \frac2{X^2+1}\Bigr) ,
\]
yaitu penguraian pada \cref{exo:b1:fractions:4}.

\textbf{11.} $P = \prod_{k=1}^{n-1}(X - \omega^k)$ (bagilah $X^n -
1$ dengan $X - 1$), jadi menurut \cref{exo:b1:poly:11}, $\frac{P'}P =
\sum_{k=1}^{n-1}\frac1{X-\omega^k}$. Di $X = 1$: $P(1) = n$ dan
$P'(1) = \sum_{j=1}^{n-1}j = \frac{n(n-1)}2$, sehingga
\[
\sum_{k=1}^{n-1}\frac1{1 - \omega^k}
= \frac{P'(1)}{P(1)} = \frac{n-1}2 .
\]

\textbf{12.} Setengah sudut: $1 - \eu^{\iu\theta} =
-2\iu\sin\frac\theta2\,\eu^{\iu\theta/2}$, sehingga
\[
\frac1{1 - \eu^{\iu\theta}}
= \frac{\eu^{-\iu\theta/2}}{-2\iu\sin\frac\theta2}
= \frac{\iu\cos\frac\theta2 + \sin\frac\theta2}
{2\sin\frac\theta2}
= \frac12 + \frac\iu2\,
\frac{\cos\frac\theta2}{\sin\frac\theta2} .
\]
Dengan $\theta = \theta_k = \frac{2k\pi}n$: $\frac1{1-\omega^k} =
\frac12 + \frac\iu2\cot\frac{k\pi}n$. Dengan menjumlahkannya atas $k = 1, \dots,
n-1$ lalu membandingkannya dengan nilai real $\frac{n-1}2$ pada
pertanyaan 11: bagian realnya sudah menanggung segalanya, sehingga
$\sum_k\cot\frac{k\pi}n = 0$ --- sebagaimana juga ditunjukkan simetri $\cot\frac{(n-k)
\pi}n = -\cot\frac{k\pi}n$.

\textbf{13.} Dengan menurunkan $\frac{P'}P = \sum_k\frac1{X -
\omega^k}$:
\[
\frac{P''}P - \Bigl(\frac{P'}P\Bigr)^2
= -\sum_{k=1}^{n-1}\frac1{(X - \omega^k)^2} .
\]
Di $X = 1$: $P''(1) = \sum_{j=2}^{n-1}j(j-1) = 2\binom n3 =
\frac{n(n-1)(n-2)}3$ (dari kesamaan tongkat hoki, atau lewat induksi), jadi
\[
\sum_{k=1}^{n-1}\frac1{(1 - \omega^k)^2}
= \Bigl(\frac{n-1}2\Bigr)^2 - \frac{(n-1)(n-2)}3
= \frac{(n-1)\bigl(3(n-1) - 4(n-2)\bigr)}{12}
= \frac{(n-1)(5-n)}{12} .
\]

\textbf{14.} Dengan mengkuadratkan rumus pertanyaan 12, dengan $c_k =
\cot\frac{k\pi}n$:
\[
\frac1{(1-\omega^k)^2}
= \Bigl(\frac12 + \frac\iu2 c_k\Bigr)^2
= \frac14 - \frac{c_k^2}4 + \frac\iu2\,c_k .
\]
Dengan menjumlahkannya dan memakai $\sum c_k = 0$ (pertanyaan 12) beserta pertanyaan 13:
$\frac{n-1}4 - \frac14\sum_k c_k^2 = \frac{(n-1)(5-n)}{12}$, sehingga
\[
\sum_{k=1}^{n-1}\cot^2\frac{k\pi}n
= (n-1) - \frac{(n-1)(5-n)}3 = \frac{(n-1)(n-2)}3 .
\]
Lalu $\frac1{\sin^2t} = 1 + \cot^2t$ memberikan
$\sum_k\frac1{\sin^2\frac{k\pi}n} = (n-1) +
\frac{(n-1)(n-2)}3 = \frac{n^2-1}3$.

\textbf{15.} Untuk $n = 3$: $\cot^2\frac\pi3 + \cot^2\frac{2\pi}3 =
\frac13 + \frac13 = \frac23 = \frac{2\cdot1}3$; dan
$\frac1{\sin^2}$ berjumlah $\frac43 + \frac43 = \frac83 =
\frac{9-1}3$. Untuk $n = 4$: $1 + 0 + 1 = 2 = \frac{3\cdot2}3$; dan $2
+ 1 + 2 = 5 = \frac{16-1}3$. Kedua rumusnya cocok.

\textbf{16.} Untuk $t \in \intoo0{\frac\pi2}$, perbandingan klasik
$\sin t < t < \tan t$ (lewat argumen luas atau kecembungan, yang
dikenal sejak sekolah menengah) memberikan, setelah kebalikannya diambil, $\cot t <
\frac1t < \frac1{\sin t}$, dengan ketiganya positif di sana; dan mengkuadratkannya
mengawetkan urutannya. Dengan $t = \frac{k\pi}n$, $1 \leq k \leq m$,
$n = 2m+1$ (sehingga $t < \frac\pi2$):
\[
\cot^2\frac{k\pi}n < \frac{n^2}{k^2\pi^2} <
\frac1{\sin^2\frac{k\pi}n} .
\]

\textbf{17.} Menurut simetri $\cot^2\frac{(n-k)\pi}n =
\cot^2\frac{k\pi}n$ dan hal serupa untuk $\sin^2$, jumlah pertanyaan 14
terparuh: $\sum_{k=1}^{m}\cot^2\frac{k\pi}n = \frac{(n-1)(n-2)}6 =
\frac{m(2m-1)}3$ dan $\sum_{k=1}^m\frac1{\sin^2\frac{k\pi}n} =
\frac{n^2-1}6 = \frac{2m(m+1)}3$. Dengan menjumlahkan pertanyaan 16 atas $k =
1, \dots, m$ lalu mengalikannya dengan $\frac{\pi^2}{n^2}$:
\[
\frac{\pi^2}{(2m+1)^2}\cdot\frac{m(2m-1)}3
\;<\; \sum_{k=1}^{m}\frac1{k^2} \;<\;
\frac{\pi^2}{(2m+1)^2}\cdot\frac{2m(m+1)}3 .
\]
Kedua batasnya menuju $\frac{\pi^2}6$ ketika $m \to \infty$ (karena rasio
$\frac{m(2m-1)}{(2m+1)^2}$ dan $\frac{2m(m+1)}{(2m+1)^2}$ sama-sama
menuju $\frac12$), sehingga menurut pengapitan jumlah parsialnya yang naik itu
konvergen dan
\[
\sum_{k=1}^{\infty}\frac1{k^2} = \frac{\pi^2}6 .
\]

\textbf{18.} Kuadratkan $\frac1{X^2-1} = \frac12\bigl(\frac1{X-1} -
\frac1{X+1}\bigr)$:
\[
\frac1{(X^2-1)^2} = \frac14\Bigl(\frac1{(X-1)^2} +
\frac1{(X+1)^2}\Bigr) - \frac12\cdot\frac1{(X-1)(X+1)} ,
\]
lalu uraikan ulang suku silangnya $\frac1{(X-1)(X+1)} =
\frac12\bigl(\frac1{X-1} - \frac1{X+1}\bigr)$ untuk memperoleh bentuk
yang dinyatakan tadi. Di $X = 0$: ruas kirinya $1$; ruas kanannya $\frac14(1 + 1) -
\frac14(-1 - 1) = \frac12 + \frac12 = 1$.

\textbf{19.} Jumlahkan pertanyaan 18 atas $n \geq 2$. Dengan $S =
\sum_{k\geq1}\frac1{k^2} = \frac{\pi^2}6$: maka $\sum_{n\geq2}
\frac1{(n-1)^2} = S$; $\sum_{n\geq2}\frac1{(n+1)^2} = S - 1 -
\frac14$; dan teleskop bercelah dua $\sum_{n\geq2}\bigl(
\frac1{n-1} - \frac1{n+1}\bigr) = 1 + \frac12 = \frac32$. Jadi
\[
\sum_{n=2}^{\infty}\frac1{(n^2-1)^2}
= \frac14\Bigl(2S - \frac54\Bigr) - \frac14\cdot\frac32
= \frac S2 - \frac{11}{16}
= \frac{\pi^2}{12} - \frac{11}{16} \approx 0.135 .
\]
Secara numerik: $\frac19 + \frac1{64} + \frac1{225} + \frac1{576} +
\frac1{1225} + \dots \approx 0.1111 + 0.0156 + 0.0044 + 0.0017 +
0.0008 + \dots \approx 0.135$: jadi sejalan.

\textbf{20.} Turunkan kesamaan pertanyaan 8
$\sum_k\frac1{X-\omega^k} = \frac{nX^{n-1}}{X^n-1}$:
\[
\sum_{k=0}^{n-1}\frac1{(X - \omega^k)^2}
= \frac{n^2X^{2n-2} - n(n-1)X^{n-2}(X^n - 1)}{(X^n - 1)^2} .
\]
Periksa di $n = 2$, $X = 2$: ruas kanannya $\frac{4\cdot4 -
2\cdot1\cdot3}{9} = \frac{10}9$; ruas kirinya $\frac1{(2-1)^2} +
\frac1{(2+1)^2} = 1 + \frac19 = \frac{10}9$.

\textbf{21.} $\frac1{\binom nk} = \frac{k!\,(n-k)!}{n!} =
\frac{k!}{(n-k+1)(n-k+2)\cdots n}$, yaitu hasil kali $k$ bilangan bulat
berurutan di penyebutnya. Dengan mensubstitusikan $j = n - k + 1$ (sehingga $j$
menjelajahi $\N^*$ ketika $n$ berjalan dari $k$):
\[
\sum_{n=k}^{\infty}\frac1{\binom nk}
= k!\sum_{j=1}^{\infty}\frac1{j(j+1)\cdots(j+k-1)}
= \frac{k!}{(k-1)\,(k-1)!} = \frac{k}{k-1} ,
\]
menurut pertanyaan 4 yang diterapkan dengan $k - 1$ menggantikan $k$ (yang sah karena
$k - 1 \geq 1$).

\textbf{22.} Untuk $k = 3$: $\frac1{\binom n3} =
\frac6{(n-2)(n-1)n}$, jadi menurut jumlah parsial pertanyaan 3 (yang tergeser),
\[
\sum_{n=3}^{N}\frac1{\binom n3}
= 6\sum_{j=1}^{N-2}\frac1{j(j+1)(j+2)}
= 6\Bigl(\frac14 - \frac1{2(N-1)N}\Bigr)
= \frac32 - \frac3{(N-1)N}
\longrightarrow \frac32 ,
\]
yaitu nilai $\frac k{k-1} = \frac32$ pada pertanyaan 21.

\textbf{23.} Untuk $n = 6$: pasangannya $k = 1$ ($\theta_1 = \frac\pi3$,
$2\cos\theta_1 = 1$) dan $k = 2$ ($\theta_2 = \frac{2\pi}3$,
$2\cos\theta_2 = -1$), ditambah \omterm{def:b1:fractions:field}{kutub} real $\pm1$:
\[
\frac1{X^6-1} = \frac16\Bigl(\frac1{X-1} - \frac1{X+1}
+ \frac{X - 2}{X^2 - X + 1}
+ \frac{-X - 2}{X^2 + X + 1}\Bigr) .
\]
Di $X = 0$: ruas kirinya $-1$; ruas kanannya $\frac16(-1 - 1 - 2 - 2)
= -1$: jadi benar.

\textbf{24.} (i) Ketunggalannya mengesahkan setiap penyamaan
koefisien --- baik penutupan, pengelompokan pasangan \omterm{def:b1:complex:field}{konjugat} pada
pertanyaan 9, maupun kiat penurunan (pertanyaan 13, 20)
semuanya bersandar padanya. (ii) \omterm{def:b1:complex:unity}{Akar satuan} memasok \omterm{def:b1:fractions:field}{kutub} $X^n -
1$, simetrinya ($k \leftrightarrow n-k$), dan
aljabar setengah sudut pada pertanyaan 12
(\cref{met:b1:complex:trig}). (iii) Adapun turunan logaritmik
$\frac{P'}P = \sum\frac{m_i}{X-a_i}$ mengubah informasi tentang
akar $P = 1 + X + \dots + X^{n-1}$ menjadi jumlah numerik
pada pertanyaan 11 dan 13 --- yaitu engsel antara Bagian II dan
Bagian III.

\textbf{25.} Penguraiannya adalah kesamaan yang murni aljabar,
benar untuk setiap nilai variabelnya sekaligus; dan itulah yang menjadikannya
sebuah mesin. Mensubstitusikan bilangan bulat lalu menjumlahkannya mengubahnya menjadi
teleskop (Bagian I); mensubstitusikan \omterm{def:b1:complex:unity}{akar satuan} lalu mengelompokkan
\omterm{def:b1:complex:field}{konjugatnya} mengubahnya menjadi kesamaan trigonometri (Bagian II dan
III); dan barulah pada langkah terakhir analisis masuk --- yaitu sebuah
pengapitan di antara dua bentuk tertutup --- untuk menyerahkan $\frac{\pi^2}6$,
yakni \omterm{def:b1:logic:statement}{pernyataan} yang tak dapat dicapai substitusi hingga mana pun. Pembagian
tugas ini (aljabar menghasilkan kesamaan hingga yang eksak, analisis melewatkannya
ke limit) adalah cetakan bagi \cref{ch:b1:series}, yang di sana
teleskop dan perbandingan menjadi sistematis, dan bagi
\cref{ch:b1:integration}, yang di sana setiap bata memperoleh antiturunan
dan penguraian yang sama menghitung integral alih-alih jumlah.

\end{solution}
